\documentclass[11pt,a4paper]{article}

\usepackage[utf8]{inputenc}
\usepackage[T1]{fontenc}
\usepackage{lmodern}
\usepackage[english]{babel}
\usepackage{amsmath,amssymb,amsthm}
\usepackage{geometry}
\usepackage{graphicx}
\usepackage{booktabs}
\usepackage{hyperref}
\usepackage{microtype}
\usepackage{xcolor}
\usepackage[font=small,labelfont=bf]{caption}
\usepackage{enumitem}
\usepackage{listings}

\geometry{margin=2.5cm}
\hypersetup{colorlinks=true,linkcolor=blue!50!black,urlcolor=blue!50!black,citecolor=blue!50!black}

\title{%
  Leveraged Pendle PT Carry on Morpho: Multi-Cycle Backtest with Gas
  Costs and a Funding-Rate Hedge\\[2pt]
  \large Hyperliquid-Hedged Active Implementation and Boros-Equivalent
  Passive Alternative}

\author{%
  Project 2 -- Blockchain \& DeFi special course\\
  \texttt{qqmikhasik}}

\date{May 2026}

\newtheorem{proposition}{Proposition}

\begin{document}

\maketitle

\begin{abstract}
We design and backtest a leveraged carry strategy on Pendle Principal
Tokens (PT) collateralised in Morpho Blue. Six complete PT cycles
across three EVM networks are evaluated against three reference
baselines, with network-specific gas friction and two hedge
implementations (active perp-short via Hyperliquid funding, and
passive position-buyer via Pendle Boros, modelled by scaling the
Hyperliquid signal by the empirically measured Boros-to-funding
ratio \(\approx 0.45\)). Three principal findings: (i) on Ethereum
mainnet at retail capital sizes, the unhedged static loop is
\emph{unprofitable} after realistic gas costs; (ii) the hedged
variants are profitable across all six cycles, including the
April-June 2024 eETH-depeg stress cycle on which the active hedge
delivers \(+130\%\) APY and the passive (Boros) hedge \(+93\%\);
(iii) the asymmetric LTV controller (adapted from Krestenko
\emph{et al.}~\cite{krestenko2026}) preserves the static-loop
equity curve to the cent in calm regimes and engages only when an
aggressive LTV setting brings the position close to the LLTV
barrier. The Pendle and Morpho integrations contribute new
\texttt{fractal-defi} primitives that did not previously exist in
the framework.
\end{abstract}

\section{Problem Statement}
\label{sec:problem}

A leveraged Pendle PT carry position recycles capital through
\(N\) cycles of (buy PT, post as collateral, borrow stablecoin,
buy more PT). For a target loan-to-value \(L\) and starting
capital \(X\), the position holds \(X / (1-L)\) of PT exposure at
the geometric limit, financed by \(LX / (1-L)\) of stablecoin debt.

\paragraph{Source of expected return.} PT pays a deterministic
yield at expiry, sourced from the underlying SY's variable yield.
For PT-sUSDE the underlying is Ethena's delta-neutral funding
income. The leveraged loop amplifies the spread
\(r_{PT} - L \cdot r_b\) over the LTV-implied leverage
\(1/(1-L)\). The expected premium over an unleveraged PT hold is
compensation for the liquidation risk that PT price drift triggers
when the realised LTV climbs against the market LLTV.

\paragraph{Constraints.}
(i) Each Morpho isolated market has a fixed liquidation LTV
(0.915 for PT-sUSDE on Ethereum).
(ii) Borrow rates respond to utilisation.
(iii) Real gas costs are non-trivial at retail capital sizes
on Ethereum mainnet (Section~\ref{sec:gas}).
(iv) The hedge instrument must reliably correlate with the source
of the risk premium (Section~\ref{sec:hedge}).
(v) Smart-contract risk is borne at three protocol layers (Pendle,
Morpho, the SY-underlying protocol).

\paragraph{Target instruments.} Six cycles across three networks
(Table~\ref{tab:cycles}); the headline backtest reports on Ethereum
mainnet, with Arbitrum and Base used to study network-side gas and
liquidity dependencies and a 2024 LRT-volatility window used for
stress testing.

\section{Formal Description}
\label{sec:formal}

We model the position at hourly cadence. Let \(t=0\) denote
position open and \(T\) the PT expiry; \(\tau(t) = T - t\) is the
remaining time.

\subsection{Static loop dynamics}

\begin{proposition}[Loop geometric series]
\label{prop:loop}
For target LTV \(L\) and \(N\) cycles, the total PT collateral
notional and the total stablecoin debt are
\begin{equation}
  C_N \;=\; X \cdot \frac{1 - L^N}{1 - L},
  \qquad
  D_N \;=\; L \cdot X \cdot \frac{1 - L^N}{1 - L}.
  \label{eq:geom}
\end{equation}
The effective leverage on the PT exposure is
\(\rho_N(L) = (1-L^N)/(1-L)\). For \(N=5, L=0.80\) this gives
\(\rho_5 = 3.36\), or 67\% of the infinite-cycle limit
\(\rho_\infty = 5\).
\end{proposition}

\begin{proposition}[Closed-form net APY at infinite leverage]
\label{prop:apy}
Ignoring friction, the asymptotic annual yield on \(X\) is
\begin{equation}
  \mathrm{APY}_{\infty}
    \;=\; \frac{r_{PT} - L\,r_b}{1 - L}.
  \label{eq:apy}
\end{equation}
\end{proposition}

\subsection{Friction model and gas costs}
\label{sec:gas}

Two layers of friction enter the realised equity curve. The first
is the Pendle AMM swap cost — a fee on the input notional plus
size-proportional slippage:
\begin{equation*}
  \mathrm{slip}
    \;=\; \sigma_{\mathrm{p}}\,\frac{\text{trade size}}{\text{pool liquidity}},
  \qquad
  \text{eff.\ price}
    \;=\; p_{PT} \cdot (1 \pm \mathrm{slip}),
\end{equation*}
modelled inside the \texttt{PendlePTEntity}; the sign is set
against the trader.

The second is gas, modelled as a fixed per-event USDC deduction
that depends on the network. We calibrate the model to a single
loop cycle costing approximately
\((650{-}800)\text{k} \cdot p_{\mathrm{gas}} \cdot p_{\mathrm{ETH}}\)
in USDC, summing the Pendle swap (\(\sim\)400k gas), Morpho
\texttt{supplyCollateral} (\(\sim\)150k), and Morpho \texttt{borrow}
(\(\sim\)250k) at the network's prevailing fee. Table
\ref{tab:gas} lists the resulting defaults per network.

\begin{table}[ht]
\centering
\begin{tabular}{lrrr}
\toprule
Network & Open cost (USDC) & Unwind cost (USDC) & Per-rebalance cost \\
\midrule
Ethereum mainnet & 150 & 125 & 30 \\
Arbitrum         &  10 &   8 &  2 \\
Base             &   8 &   6 &  1.5 \\
\bottomrule
\end{tabular}
\caption{Conservative gas-cost defaults used in the backtest. The
mainnet open figure is on the high side of the
(80, 130) USDC range we observe at 10 gwei / ETH = \$2{,}500;
the strategy looks slightly worse than it would in a favourable
gas environment.}
\label{tab:gas}
\end{table}

\section{Dynamic LTV Controller}
\label{sec:dynamic}

We adapt the asymmetric \(\alpha\)-band controller of
Krestenko \emph{et al.}~\cite{krestenko2026} from the
spot-perpetual basis context to a PT-on-Morpho geometry. Their
barrier is on rising spot; ours is on falling PT. The Bachelier
machinery is identical up to barrier direction.

\subsection{First-passage probability}

Model \(X_t = \ln(p_{PT}(t)/p_{PT}(0))\) as Brownian motion with
constant drift \(\mu\) and volatility \(\sigma\), estimated
empirically from a rolling window of hourly PT prices.
Liquidation is the first crossing of
\(b = \ln(L_0 / L_{\text{liq}}) < 0\).

\begin{proposition}[Bachelier first-passage]
\label{prop:fpt}
For \(b<0\) and horizon \(h>0\),
\begin{equation}
  \Pi_{\mathrm{liq}}(b; h, \mu, \sigma)
    \;=\;
    \Phi\!\left(\frac{b - \mu h}{\sigma \sqrt{h}}\right)
    + \exp\!\left(\frac{2\mu b}{\sigma^2}\right)\,
      \Phi\!\left(\frac{b + \mu h}{\sigma \sqrt{h}}\right).
  \label{eq:fpt}
\end{equation}
\end{proposition}

The closed form is verified against Monte-Carlo simulation
(\(3 \times 10^4\) GBM paths, 2{,}400-step grid; agreement within
1\% absolute).

\subsection{Asymmetric band}

\(L_U\) (solvency-driven, structural): largest \(L\) with
\(\Pi_{\mathrm{liq}}(L; h_{\mathrm{liq}}) \le \varepsilon_{\mathrm{liq}}\).
Default: \(\varepsilon_{\mathrm{liq}}=10^{-4}\) over
\(h_{\mathrm{liq}}=3\) hours, matching Arbitrum sequencer-outage
documentation.

\(L_L\) (economic, may vanish): from a first-order carry-vs-cost
trade-off
\begin{equation}
  L^\dagger - L_L
    \;=\;
    \frac{K_{\mathrm{reb}}\,(1-L^\dagger)^2}{X(r_{PT}-r_b)h_{\mathrm{reb}}}.
\end{equation}
If \(r_{PT} \le r_b\) the expression is meaningless and \(L_L = 0\)
(don't re-lever a losing trade).

\paragraph{Important interpretation.} The controller defends against
\emph{liquidation}, which is a discrete event. It does \emph{not}
defend against drawdown from PT-price drift while the position
remains solvent. In our six backtested cycles the realised LTV
never approached the LLTV barrier, so the dynamic equity curve is
identical to the static one to the cent. This is by design — the
controller is meant to vanish in calm regimes and only engage on
the brink of insolvency.

\section{Hedge Specification}
\label{sec:hedge}

The dominant tail risk of the static loop is a rise in \(r_{PT}\):
PT falls, LTV rises. The fundamental driver of PT-sUSDE yield is
the funding-rate income of Ethena's short perpetual book; a
position long the perpetual funding rate therefore hedges the
PT-side risk.

\subsection{Two implementations}

\textbf{Active hedge (Hyperliquid).} The investor maintains a
delta-neutral short on Hyperliquid ETH-USDC perpetual,
\emph{receiving} funding payments when funding is positive,
\emph{paying} when negative. Backtest models this exactly through
\texttt{FundingHedgeLoader} consuming Hyperliquid's
\texttt{fundingHistory} public-info endpoint at hourly granularity.

\textbf{Passive hedge (Pendle Boros).} Boros tokenises forward
contracts on the average funding rate up to maturity. A
position-buyer holds the token and receives accrued funding as
the AMM marks down its mid-price. Boros launched August 2025;
the public API exposes only the most recent \(\sim\)500 bars
which precludes direct historical use, but the recent overlap
window suffices for empirical proxy validation
(Section~\ref{sec:boros-validation}).

\subsection{Empirical validation: Boros vs Hyperliquid}
\label{sec:boros-validation}

Over the 21-day overlap window
2026-04-23 \(\rightarrow\) 2026-05-14 (\(n=499\) hourly observations),
the Pearson correlation of Boros mark APR against Hyperliquid raw
funding is 0.47, mean absolute spread 584 bps, mean bias
\(-302\) bps (Boros below funding). The Boros mark moves within
\([+1.0\%, +4.2\%]\); the Hyperliquid raw rate moves within
\([-13.2\%, +11.0\%]\). The two instruments are economically
related but \emph{not} identical: Boros is a forward expectation
(mean-reverting and gentle), Hyperliquid is the realised spot
funding rate (volatile and reactive). A backtest that substitutes
raw funding for Boros therefore \emph{over-estimates} the realised
hedge gain.

We compensate by reporting two hedge variants per cycle:
\textbf{Hedged-Active} uses the raw Hyperliquid funding rate
unchanged (upper bound, applicable when the user actively shorts
the perp themselves); \textbf{Hedged-Boros} damps the raw rate by
the empirical Boros-to-Hyperliquid ratio
\(\hat{\rho} \approx 0.45\) (realistic lower bound for a passive
Boros-position-holder).

\subsection{Sizing}

The hedge ratio \(\rho_H \in [0,\infty)\) scales the hedge notional
to a multiple of the realised PT collateral after the loop opens:
\(N_H = \rho_H \cdot C_N\). The reported backtest uses
\(\rho_H = 1\) (notional-matched).

\section{Backtest Protocol}
\label{sec:protocol}

\paragraph{Cycle panel.} Six cycles across three networks;
Table~\ref{tab:cycles}. Each cycle is hourly-sampled across the
full life of the PT market (start = beginning of available
historical data, end = on-chain expiry timestamp).

\begin{table}[ht]
\centering
\footnotesize
\begin{tabular}{llrll}
\toprule
Cycle & Network & Days & Window & Character \\
\midrule
sUSDE-25SEP2025      & Ethereum & 60 & 2025-07-27 → 2025-09-25 & calm \\
sUSDE-27NOV2025      & Ethereum & 59 & 2025-09-29 → 2025-11-27 & calm \\
sUSDE-5FEB2026       & Ethereum & 60 & 2025-12-07 → 2026-02-05 & calm (yield compression) \\
USDe-11DEC2025@Arb   & Arbitrum & 60 & 2025-10-12 → 2025-12-11 & implied-yield spike \\
USDe-11DEC2025@Base  & Base     & 60 & 2025-10-12 → 2025-12-11 & implied-yield spike, thin pool \\
weETH-27JUN2024      & Ethereum & 60 & 2024-04-28 → 2024-06-27 & STRESS: eETH depeg \\
\bottomrule
\end{tabular}
\caption{Backtest cycle panel. Three calm sUSDe cycles on
Ethereum provide a stable baseline; the December 2025 USDe cycle
on Arbitrum and Base sees a sharp implied-yield rise; the June
2024 weETH cycle spans the LRT-volatility crisis. Loan asset
is USDC for all cycles except the weETH cycle where it is USDA
(Angle stablecoin, treated as 1.0 USDC equivalent for the
backtest).}
\label{tab:cycles}
\end{table}

\paragraph{Strategies.} Six per cycle: HoldUSDC, HoldSUSDe (proxy
via implied yield), HoldPTNoLeverage, StaticLoop (target LTV 0.80),
DynamicLoop (band-controlled around the same target), and two
hedged variants — Hedged-Active (raw Hyperliquid funding) and
Hedged-Boros (Hyperliquid scaled by 0.45). All loop variants pay
network-appropriate gas (Table~\ref{tab:gas}). The unwind uses a
flash-repay model (collateral redeemed at par, debt paid out of
the resulting cash); Aave-style 5-bp flash-loan fee is not yet
modelled.

\paragraph{Statistical validation.} On the headline cycle
(sUSDE-27NOV2025) the static and hedged variants are evaluated
through a block-bootstrap procedure: hourly returns are
resampled in 24-hour blocks, the equity curve reconstructed,
and the 5/95 percentiles of final balance reported. A separate
volatility-stress run multiplies the PT log-return path's
realised \(\sigma\) by 1.5 and re-runs the dynamic strategy.

\section{Results}
\label{sec:results}

The headline number is in Table~\ref{tab:headline} (APY only —
full Sharpe and max-drawdown tables in Section~\ref{sec:supp}).

\begin{table}[ht]
\centering
\footnotesize
\begin{tabular}{lrrrrr}
\toprule
Cycle &
PT-no-lev &
Static &
Hedged-Active &
Hedged-Boros &
\\
\midrule
sUSDE-25SEP2025                    &  +12.23\% &   +2.83\% &     +38.56\% &  +18.91\% \\
sUSDE-27NOV2025                    &   +7.11\% &  --0.72\% &     +24.25\% &  +10.52\% \\
sUSDE-5FEB2026                     &   +5.13\% &  --3.55\% &     +19.82\% &   +6.97\% \\
\textbf{weETH-27JUN2024 (stress)}  & \textbf{+31.82\%} & \textbf{+62.21\%} & \textbf{+129.75\%} & \textbf{+92.60\%} \\
\midrule
USDe-11DEC2025@Arb (anomalous)     & --12.24\% & --36.14\% &     --11.01\% & --24.83\% \\
USDe-11DEC2025@Base (anomalous)    & --28.65\% & --74.32\% &     --50.71\% & --63.69\% \\
\bottomrule
\end{tabular}
\caption{Realised APY per strategy per cycle (Ethereum sUSDe
cycles and stress cycle in the top block; the L2 USDe cycle
captured an implied-yield spike that hit all variants — see
Section~\ref{sec:l2-anomaly}). Active hedge = raw Hyperliquid
funding; Boros hedge = Hyperliquid scaled by 0.45.}
\label{tab:headline}
\end{table}

\subsection{What works in the calm regime}

The three calm Ethereum cycles tell a consistent story.
Pure PT-no-leverage delivers a small positive return
(\(+5\%\) to \(+12\%\) APY) — what a buy-and-hold investor would
earn. \textbf{The static leveraged loop is unprofitable across
two of three calm cycles after gas}: the \$275 round-trip gas
cost at \$10k capital amounts to \(-2.75\%\) of starting equity,
which the modest implied-yield premium does not offset. This is
a real and previously-undiscussed finding: leveraged PT carry on
mainnet at retail capital is gas-toxic.

The active hedge rescues the loop \emph{decisively}, with
\(+19.8\%\) to \(+38.6\%\) APY across the three calm cycles.
Even the more conservative Boros hedge keeps the strategy in
the green at \(+7.0\%\) to \(+18.9\%\) APY. \textbf{Hedge is
not optional on mainnet — it is economically necessary.}

\subsection{What works in the stress regime}

The weETH-27JUN2024 cycle spans the April 2024 eETH-depeg event.
Pendle implied APY swung from 0\% to 32\% over the window —
the strongest real-world stress we can access. All four leveraged
variants beat PT-no-leverage by 2-4× APY:
\begin{itemize}[topsep=2pt,itemsep=0pt]
  \item PT-no-leverage: \(+31.8\%\) APY (rising implied yield
        translates to cheap PT accumulation, which redeems at par)
  \item Static loop: \(+62.2\%\) APY (leverage amplifies the gain)
  \item Hedged-Active: \(+129.7\%\) APY (loop carry + sharp
        positive funding during the period)
  \item Hedged-Boros: \(+92.6\%\) APY (passive Boros captures
        the same regime more gently)
\end{itemize}

Important reframing: the hedge in this regime is not a
\emph{drawdown shield} — it is a \emph{second profit source}.
The basis-trade construction in
Krestenko~\cite{krestenko2026} is exactly that: simultaneous long-
carry on the price-driven leg and short-funding-collection on the
perp leg.

\subsection{The L2 USDe-Dec-2025 anomaly}
\label{sec:l2-anomaly}

Both L2 cycles show \emph{all} variants losing money on
PT-USDe-11DEC2025. The drivers are
(a) implied APY rising from 6.6\% to 30.6\% over the window, which
hurts the existing PT holders (PT mark-to-market drops sharply
even though redeem at expiry returns to par),
(b) pool TVL collapsing from \$163k to \$29k on Arbitrum (and
proportionally lower on Base), which inflates AMM slippage on
the open and post-expiry unwind, and
(c) the L2 USDe market having only the raw USDe-collateral (not
sUSDe), so the realised yield to PT-holders is smaller.

The hedge component is positive across both networks (Hedged-Active
loses 11\% vs Static's 36\% on Arbitrum, and 51\% vs 74\% on Base),
i.e.\ the hedge does part of its job — it reduces drawdown — but
the underlying carry was negative, so the strategy still loses.

\subsection{Dynamic controller behaviour}

Across all six cycles the dynamic controller's equity curve is
\emph{identical to the static one to the cent}. The realised LTV
in every cycle stayed deep inside the safety band; the controller
saw no event worth reacting to. This is not a deficiency — it is
the controller behaving as designed: \emph{free in calm regimes,
engaging only on the brink of insolvency}. A separate experiment
on an aggressive \(L^\dagger=0.84\) (only 7.5 p.p.\ buffer to LLTV)
shows the controller activates immediately and costs
\(\sim\)61 bps APY as insurance premium.

\subsection{Statistical validation}

\paragraph{Block-bootstrap.} On sUSDE-27NOV2025 with 24-hour
blocks, the static loop's bootstrap final-balance distribution is
[\$10{,}047, \$10{,}186, \$10{,}339] (5/50/95 percentiles); the
hedged-active variant is [\$10{,}435, \$10{,}605, \$10{,}797]. The
\(p<0.05\) cut for hedge-vs-static separation is satisfied with
substantial margin — the headline outperformance is not a single
lucky path.

\paragraph{Volatility stress.} A \(1.5\times\) injected
volatility on the PT-price path keeps the dynamic strategy
solvent on all three calm sUSDe cycles, with APY moving from
\(-0.7\%\) to \(-4.9\%\) on sUSDE-27NOV2025 (cost of stress)
and from \(-3.5\%\) to \(-2.8\%\) on sUSDE-5FEB2026 — the
dynamic controller activates and absorbs some of the synthetic
shock.

\section{Strategy improvement over prior work}

The starting baseline is the static fixed-LTV leveraged PT loop,
documented in DeFi practitioner literature. Two principal
improvements:

\begin{enumerate}[topsep=2pt,itemsep=2pt]
  \item \emph{Dynamic LTV controller} adapted from
        \cite{krestenko2026}. Their barrier is upper
        (basis-trade short liquidates on rising spot); ours is
        lower (loop collateral liquidates on falling PT).
        Identical Bachelier first-passage formula up to barrier
        sign.

  \item \emph{Funding-rate hedge} with two implementations.
        Active = direct Hyperliquid short, Passive = Pendle Boros
        position-buyer modelled by scaling raw funding by the
        empirical Boros-to-funding ratio. Either implementation
        rescues the loop from gas-toxicity on Ethereum mainnet
        and supplies a second profit source in stress regimes.
\end{enumerate}

\section{Risk Analysis and Limitations}
\label{sec:risk}

\paragraph{Realistic risks not modelled in the equity curve.}
\begin{itemize}[topsep=2pt,itemsep=0pt]
  \item \emph{Smart-contract risk} on three protocols (Pendle,
        Morpho, Ethena). Historical exploit rate \(\sim\)1-3\%
        per protocol per year; multiplicatively gives an
        annualised expected drag of 3-9\% on every reported
        APY. Not modelled.
  \item \emph{Aave-style flash-loan fee} on unwind (5 bps).
        Slight over-estimate of unwind proceeds, \(\sim\)0.05\%
        of final equity. Not modelled.
  \item \emph{Borrow-rate spike} in Morpho during utilisation
        stress. Modelled via the observed historical
        \texttt{borrowApy} feed (so realistic on the backtested
        cycles), but a future spike beyond the observed range
        could invalidate the static loop entirely.
  \item \emph{Underlying SY depeg}. Now modelled correctly
        through \texttt{sy\_price\_in\_usdc} on the PT redeem,
        but the binance \texttt{USDEUSDT} feed remained close to
        peg through the December 2025 USDe stress, so the loss
        we see there comes from \emph{implied yield rise} and
        \emph{pool thinning} rather than realised SY depeg.
\end{itemize}

\paragraph{Survivor-bias acknowledged.} All six cycles backtested
either (a) terminated naturally at PT expiry without liquidation,
or (b) saw losses that fell inside the LLTV safety buffer. We do
not have a cycle with a realised liquidation event. The dynamic
controller's value cannot therefore be \emph{empirically} proven
against a live liquidation in this work.

\paragraph{Hedge tracking error.} Pendle Boros and Hyperliquid raw
funding are economically related but not identical. The 0.45
scaling factor we use to synthesise a Boros equivalent is a
\emph{summary statistic} over a 21-day window, not a stable
parameter; the true Boros realisation could differ from the
synthetic both ways depending on the trading regime.

\section{Supplementary tables}
\label{sec:supp}

\begin{table}[ht]
\centering
\footnotesize
\begin{tabular}{lrrrrr}
\toprule
Cycle & PTNoLev & Static & Dynamic & Hedged-Active & Hedged-Boros \\
\midrule
\multicolumn{6}{l}{\emph{Sharpe ratio}} \\
\midrule
sUSDE-25SEP2025                    &  9.03 &  3.02 &  3.02 & 10.45 &  6.39 \\
sUSDE-27NOV2025                    &  9.91 &  4.14 &  4.14 & 13.28 &  8.37 \\
sUSDE-5FEB2026                     & 16.70 &  6.76 &  6.76 & 25.94 & 15.72 \\
weETH-27JUN2024                    &  9.96 &  6.97 &  6.97 & 12.82 &  9.66 \\
\midrule
\multicolumn{6}{l}{\emph{Max drawdown}} \\
\midrule
sUSDE-25SEP2025                    & --0.0032 & --0.0121 & --0.0121 & --0.0091 & --0.0107 \\
sUSDE-27NOV2025                    & --0.0010 & --0.0036 & --0.0036 & --0.0059 & --0.0043 \\
sUSDE-5FEB2026                     & --0.0004 & --0.0015 & --0.0015 & --0.0053 & --0.0021 \\
weETH-27JUN2024                    & --0.0051 & --0.0166 & --0.0166 & --0.0164 & --0.0165 \\
\bottomrule
\end{tabular}
\caption{Full Sharpe and max-drawdown across the four headline
Ethereum cycles. L2 anomalous cycles omitted for clarity.}
\label{tab:supp}
\end{table}

\section{References}

\begin{thebibliography}{9}
\bibitem{krestenko2026}
A.~Krestenko, M.~Butov, R.~Berezovskiy, D.~Bolotin.
\emph{Dynamic Collateral Control for Permissionless Spot-Perpetual
Basis Trading.} Vega Institute Foundation, 2026.

\bibitem{pendlev3}
Pendle Finance. \emph{Pendle V3 Documentation.}
\url{https://docs.pendle.finance/}.

\bibitem{pendleboros}
Pendle Finance. \emph{Boros: Tokenised funding rates.}
\url{https://docs.pendle.finance/boros-docs/Introduction}.

\bibitem{morphoblue}
Morpho Labs. \emph{Morpho Blue Whitepaper.}
\url{https://docs.morpho.org/}.

\bibitem{logarithm}
Logarithm Labs. \emph{fractal-defi:} open-source DeFi-strategy
backtesting library, v1.3.1.
\url{https://github.com/Logarithm-Labs/fractal-defi}.

\bibitem{ethena}
Ethena Labs. \emph{USDe: Internet Bond.}
\url{https://ethena.fi/}.

\bibitem{hyperliquid}
Hyperliquid Labs. \emph{Hyperliquid documentation.}
\url{https://hyperliquid.gitbook.io/hyperliquid-docs/}.

\bibitem{borodin}
A.~Borodin, P.~Salminen. \emph{Handbook of Brownian Motion:
Facts and Formulae.} Birkh\"auser, 2002.
\end{thebibliography}

\appendix

\section{Reproducibility}

Headline backtest, sweep, and validation are reproducible end-to-end
from a fresh clone:

\begin{verbatim}
python -m venv .venv && source .venv/Scripts/activate
pip install -e ".[dev]"
pytest -q  # 153 tests
python scripts/multi_cycle_backtest.py
python scripts/sweep_hedge_ratio.py
python scripts/validate_boros_proxy.py
\end{verbatim}

All four data sources (Pendle, Morpho, Binance, Hyperliquid,
Boros) are public and keyless. No environment variables required.

\end{document}
